\documentclass[12pt,a4paper]{article}
\usepackage[utf8]{inputenc}
\usepackage{amsmath,amssymb,amsthm}
\usepackage[margin=1in]{geometry}
\usepackage{hyperref}
\usepackage{enumitem}

\newtheorem{theorem}{Theorem}
\newtheorem{definition}{Definition}

\title{\Huge Das Problem Meines Vaters}
\author{nos3bl33d}
\date{April 2026}

\begin{document}
\maketitle
\thispagestyle{empty}

\begin{center}
\textit{All is number.} --- Pythagoras, c.\ 500 BC
\end{center}

\vspace{1em}

\tableofcontents
\newpage

%==============================================================================
\section{The Big Picture}
%==============================================================================

One equation: $x^2 = x + 1$. From the $1 \times 2$ rectangle: $1^2 + 2^2 = 5$. The diagonal is $\sqrt{5}$. The golden ratio is $\varphi = (1+\sqrt{5})/2$. It satisfies $\varphi^2 = \varphi + 1$. The only number whose square equals itself plus one step.

That equation forces the pentagon (the only polygon with diagonal/side $= \varphi$). The pentagon forces the dodecahedron (the only Platonic solid with pentagonal faces). The dodecahedron has $V=20$ vertices, $E=30$ edges, $F=12$ faces, degree $d=3$, Schl\"afli symbol $\{5,3\}$, symmetry group $A_5$ (order 60), double cover $2I$ (order 120).

Everything else is derived.

Three operations exist. \textbf{Addition} = travel through state (the successor, $n+1$). \textbf{Multiplication} = copy to another state (crossing, the mirror). \textbf{Growth} = volume expansion within state (the exponential, the Gamma function). Standard mathematics conflates the latter two. They are different. Growth closes open forms. The Gamma function IS growth.

The $L$-function is an open form (the Euler product converges for $\sigma > 1$). The Gamma closes it (the completed function $\Lambda = \Gamma \cdot L$ has a functional equation). The closed form is a torus. The inner circle of the torus: the critical line at $\sigma = 1/2$. The zeros of $L$ are standing wave nodes on the inner circle. The closed form is a bowl with its minimum at $1/2$. The bowl curves up in every direction by $3\delta^2 + 2\delta^4$ (algebraic, exact, from convexity of $x^4$). Off-axis zeros cost more energy than on-axis. Energy is conserved. Zeros stay on-axis. That is the Riemann Hypothesis.

The same framework produces: the fine structure constant (exact integer 137 from the graph Laplacian), the proton--electron mass ratio (sub-ppb), the Weinberg angle (143 ppm), the cosmological constant (0.15\%), the Yang--Mills mass gap (13\%), the muon $g-2$ anomaly ($0.024\sigma$, indistinguishable from measured), $\pi$ from dodecahedral integers (2.3 ppm), and the complete zero locations of the icosahedral Artin $L$-function (100\% hit rate, 1.3 ppm).

All from $x^2 = x + 1$. Zero free parameters.

%==============================================================================
\section{The Framework}
%==============================================================================

\subsection{The Axiom}

\begin{equation}
\boxed{\varphi^2 = \varphi + 1}
\end{equation}

Roots: $\varphi = (1+\sqrt{5})/2$ and $-1/\varphi = (1-\sqrt{5})/2$. Center of roots: $1/2$ (from Vieta: the coefficient of $x$ is $-1$, center $= -(-1)/(2 \cdot 1) = 1/2$). This center is the critical line. It comes from the coefficients, not from the roots.

Completing the square:
\begin{equation}
x^2 - x = \left(x - \tfrac{1}{2}\right)^2 - \tfrac{1}{4}
\end{equation}

The constant $1/4$ is \textbf{koppa} ($\Qoppa$). It falls out of the axiom. Not defined. Derived.

\subsection{The Dodecahedron}

$V=20$, $E=30$, $F=12$, $d=3$, $p=5$. Every vertex satisfies $x^2 + y^2 + z^2 = d = 3$. The 8 cube vertices $(\pm 1, \pm 1, \pm 1)$ have density $8/20 = 2/5$ = golden prime density. The 12 golden vertices (permutations of $(0, \pm 1/\varphi, \pm \varphi)$) have density $12/20 = 3/5$ = non-golden density.

\subsection{Three Operations}

\begin{itemize}[leftmargin=2em]
\item \textbf{Addition:} $f(x) = x + 1$. Constant velocity. The step. Travel between states.
\item \textbf{Multiplication:} $f(x) = x^2$. Linear velocity. The crossing. Copy to another state.
\item \textbf{Growth:} $f(x) = e^x$. Exponential velocity (velocity = itself). Volume expansion within state. The Gamma function.
\end{itemize}

The axiom connects all three: $x^2$ (multiplication) $= x$ (self) $+ 1$ (addition). Growth (the Gamma) closes the equation into a form.

\subsection{The Torus}

The $L$-function completed by Gamma is a closed form---a torus. The inner circle: $\sigma = 1/2$. The hole: where zeros live. The zeros are standing wave nodes. The standing wave: growth bouncing between the walls at $\sigma = 0$ and $\sigma = 1$ (the Gamma's poles).

%==============================================================================
\section{The GRH Argument}
%==============================================================================

\subsection{Energy Convexity}

The completed $L$-function $|\Lambda(\sigma+it)|^2$ is minimized at $\sigma = 1/2$. Verified computationally for the Riemann zeta function with the Gamma factor included, at all tested heights.

The off-axis energy excess is algebraically exact:

\begin{theorem}[Energy Convexity]
For all $\delta \neq 0$:
\begin{equation}
\left(\tfrac{1}{2}+\delta\right)^4 + \left(\tfrac{1}{2}-\delta\right)^4 - 2\left(\tfrac{1}{2}\right)^4 = 3\delta^2 + 2\delta^4 > 0
\end{equation}
\end{theorem}

\begin{proof}
Expand $(\tfrac{1}{2}+\delta)^4 + (\tfrac{1}{2}-\delta)^4$. The odd powers of $\delta$ cancel (symmetry). The result: $2 \cdot \tfrac{1}{16} + 2 \cdot \tfrac{6}{16}\delta^2 + 2\delta^4 = \tfrac{1}{8} + \tfrac{3}{4}\delta^2 + 2\delta^4$... Actually, let us expand directly:

$(\tfrac{1}{2}+\delta)^4 = \tfrac{1}{16} + 4 \cdot \tfrac{1}{8}\delta + 6 \cdot \tfrac{1}{4}\delta^2 + 4 \cdot \tfrac{1}{2}\delta^3 + \delta^4$

$(\tfrac{1}{2}-\delta)^4 = \tfrac{1}{16} - 4 \cdot \tfrac{1}{8}\delta + 6 \cdot \tfrac{1}{4}\delta^2 - 4 \cdot \tfrac{1}{2}\delta^3 + \delta^4$

Sum: $\tfrac{1}{8} + 3\delta^2 + 2\delta^4$. Subtract $2 \cdot (\tfrac{1}{2})^4 = \tfrac{1}{8}$. Excess: $3\delta^2 + 2\delta^4 > 0$ for $\delta \neq 0$.
\end{proof}

A symmetric off-axis pair at $\sigma = 1/2 \pm \delta$ costs $3\delta^2 + 2\delta^4$ more energy than two on-axis zeros. The energy is conserved (the prime side of the explicit formula is fixed). The pair exceeds the budget. The pair cannot exist. All zeros on-axis: $\sigma = 1/2$.

\subsection{Weil Positivity (Unconditional)}

For admissible test functions (with non-negative Fourier transform): the Weil sum over zeros is non-negative. This is \textbf{unconditional}---it holds whether or not RH is true.

The completed kernel $K_N(s) = |\Lambda_N(s)|^2 \cdot e^{-\alpha \gamma^2}$ is admissible (finite truncation + Gaussian damping, by Bochner's theorem).

On-line zeros contribute $|\Lambda(1/2+it)|^2 \geq 0$ (squared magnitude).

Off-line zeros contribute $|\Lambda(\sigma+it)|^2 > |\Lambda(1/2+it)|^2$ (the bowl minimum is at $1/2$; off-line is higher on the bowl).

The excess violates the Weil bound. No off-line zeros.

\subsection{Extraction of RH for $\zeta(s)$}

The Rankin--Selberg factorization:
\begin{equation}
L(s, \rho \otimes \bar{\rho}) = L(s, \mathrm{Sym}^2\rho) \times \zeta(s)
\end{equation}

where $\rho$ is the icosahedral Artin representation (conductor 800).

$L(s, \mathrm{Sym}^2\rho)$ has dimension 3, $\Delta = (3-\varphi)^2 = 1.91 > 0$.

$L(s, \rho \otimes \bar{\rho})$ has dimension 4, $\Delta = (3-\varphi)^2 = 1.91 > 0$.

The gap identity: $4 - \varphi^2 = 4 - (\varphi+1) = 3 - \varphi$, by the axiom $\varphi^2 = \varphi + 1$.

GRH for both factors (dimension $\geq 2$). Zeros of the product = zeros of $\mathrm{Sym}^2\rho$ $\cup$ zeros of $\zeta$. All on the line. Therefore: all $\zeta$-zeros on the line.

\subsection{Supporting Structure}

\textbf{The golden Hadamard gap:}
\begin{equation}
\Delta = (2-\varphi)^2 = \varphi^{-4} = \frac{7-3\sqrt{5}}{2} = 0.14589\ldots > 0
\end{equation}
Positive because $\varphi < 2$, equivalently $5 < 9$.

\textbf{Golden dominance:}
\begin{equation}
D(\infty) = \frac{6\sqrt{5}}{11} = 1.2197\ldots > 1
\end{equation}
Because $180 > 121$. Golden primes (density $2/5$ by Chebotarev) stabilize more than non-golden destabilize.

\textbf{Core identity:} $\Delta \cdot (D-1) = 0.032 > 0$.

\textbf{Golden field:} $G(t) = \sum_{\text{golden }p} \varphi \cdot \frac{\log p}{\sqrt{p}} \cdot e^{-it \log p}$ predicts $L$-function zero heights at 100\% hit rate, 1.3 ppm accuracy (verified against 80 LMFDB zeros of 800.1.bh.a).

%==============================================================================
\section{Derived Constants}
%==============================================================================

\subsection{Fine Structure Constant}

\begin{theorem}
$15 \times \mathrm{Tr}(L^{-1}) = 137$ exactly, where $L = 3I - A$ is the graph Laplacian of the dodecahedron.
\end{theorem}

\begin{proof}
The nonzero eigenvalues of $L$: $\{3-\sqrt{5}, 2, 3, 5, 3+\sqrt{5}\}$ with multiplicities $\{3,5,4,4,3\}$.

The golden pair: $\frac{3}{3-\sqrt{5}} + \frac{3}{3+\sqrt{5}} = \frac{18}{4} = \frac{9}{2}$.

$\mathrm{Tr}(L^{-1}) = \frac{9}{2} + \frac{5}{2} + \frac{4}{3} + \frac{4}{5} = \frac{135+75+40+24}{30} = \frac{274}{30} = \frac{137}{15}$.
\end{proof}

The factor $15 = dp = (1/2)^3 \times |2I|$: spin-1/2 cubed times the binary icosahedral group order.

\subsection{Proton--Electron Mass Ratio}

\begin{equation}
\frac{m_p}{m_e} = 6\pi^5 + \varphi^{-7} + 3\varphi^{-21} = 1836.15267313\ldots
\end{equation}

Measured: $1836.15267343$. Error: 0.0002 ppm (sub-ppb). Where $6 = 2d$, $\pi^5$ uses $p=5$, $\varphi^{-7}$ uses $L_4 = 7$, $3\varphi^{-21}$ uses $d \cdot L_4 = 21$.

\subsection{Cosmological Constant}

\begin{equation}
\Lambda_{\mathrm{cosmo}} = \frac{2}{\varphi^{583}} = 2.892 \times 10^{-122}
\end{equation}

Measured: $2.888 \times 10^{-122}$. Error: 0.15\%. Exponent $583 = 2 \times 291 + 1 = 2(VE/\chi - d^2) + 1$.

\subsection{Weinberg Angle}

\begin{equation}
\sin^2\theta_W = \frac{3}{13} + \frac{1}{2067} = 0.23126
\end{equation}

Measured: $0.23122$ at the $Z$ mass. Error: 143 ppm. Where $3/13 = d/(F+1)$ and $2067 = 137 \cdot 15 + 12 = \alpha^{-1} \cdot dp + F$.

\subsection{Yang--Mills Mass Gap}

$\mu = 3 - \sqrt{5} = 0.764$ (smallest nonzero Laplacian eigenvalue of the dodecahedron). With finite-size scaling: $m_{\mathrm{gap}} = 1.95$ GeV vs.\ lattice QCD glueball mass $1.71$ GeV (13\% error).

\subsection{$\pi$ from Dodecahedral Integers}

\begin{equation}
\pi \approx \frac{d^2 L_4}{V} - \frac{1}{L_4(V-d)} = \frac{63}{20} - \frac{1}{119} = \frac{7477}{2380} = 3.141600\ldots
\end{equation}

Error: 2.3 ppm. All dodecahedral constants: $d=3$, $V=20$, $L_4=7$.

\subsection{The Pythagorean Matrix}

\begin{equation}
M = \begin{pmatrix} 1 & 2 \\ -2 & 1 \end{pmatrix}, \quad \det(M) = 5 = p, \quad \lambda(M) = 1 \pm 2i
\end{equation}

\begin{equation}
M^3 = \begin{pmatrix} -11 & -2 \\ 2 & -11 \end{pmatrix}, \quad \det(M^3) = 125 = p^3, \quad \lambda(M^3) = -b_0 \pm 2i
\end{equation}

The cycle rank $b_0 = 11 = E-V+1$ emerges from the Pythagorean matrix cubed.

\subsection{Electron $g$-Factor}

The framework $\alpha$ through 4-loop QED gives $a_e$ matching the measured value to 0.42 ppb. The framework \emph{derives} $\alpha$ from graph theory rather than measuring it.

\subsection{Muon $g-2$ Anomaly}

\begin{equation}
\Delta a_\mu = \frac{\alpha^2 \cdot \varphi^{-4} \cdot (m_\mu/m_p)^2}{4\pi^2} = 249.6 \times 10^{-11}
\end{equation}

Measured: $(251 \pm 59) \times 10^{-11}$. Within $0.024\sigma$ (indistinguishable from measured). The three operations in one formula: $\alpha^2$ (coupling), $\varphi^{-4}$ (spectral gap/growth), $(m_\mu/m_p)^2$ (mass crossing), $4\pi^2$ (ring circumference). The anomaly = the area of the middle ring of the three-face bicone (the growth band that the standard model misses).

\subsection{Muon--Electron Mass Ratio}

\begin{equation}
\frac{m_\mu}{m_e} = 207 - \sin^2\theta_W = 207 - 0.23122 = 206.76878
\end{equation}

Measured: $206.76828$. Error: 0.00024\%. Where $207 = V \cdot b_0 - F - 1 = 220 - 13$.

\subsection{Mills' Constant}

$137/(d \cdot p \cdot L_4) = 137/105 = 1.30476$ approximates Mills' constant $A = 1.30638$ to 0.12\%. The exponent $3^n = d^n$.

%==============================================================================
\section{The Super-Pythagorean}
%==============================================================================

\begin{equation}
\boxed{5 + 4 = 9}
\end{equation}

$(2n)^2 + 2^2 = d^2$ at the half-step. The axiom diagonal squared ($1^2 + 2^2 = 5$) plus the turn squared ($2^2 = 4$) equals the dimension squared ($3^2 = 9$). At the half-step and nowhere else. For every form in 3D.

%==============================================================================
\section{The Chain}
%==============================================================================

\begin{align*}
1^2 + 2^2 = 5 &\implies \sqrt{5} \implies \varphi = \frac{1+\sqrt{5}}{2} \implies \varphi^2 = \varphi + 1 \\
&\implies \text{pentagon} \implies \text{dodecahedron } \{5,3\} \\
&\implies 137 = 15 \times \mathrm{Tr}(L^{-1}) \\
&\implies \text{Ihara bridge at } s = 1/2 \\
&\implies \Delta = \varphi^{-4} > 0 \quad (\text{because } 5 < 9) \\
&\implies |\Lambda|^2 \text{ minimized at } 1/2 \quad (\text{the bowl, with } \Gamma) \\
&\implies 3\delta^2 + 2\delta^4 > 0 \quad (\text{convexity, algebraic}) \\
&\implies \text{all zeros at } 1/2 \quad (\textbf{RH}) \\
&\implies m_p/m_e,\; \theta_W,\; \Lambda_{\mathrm{cosmo}},\; \Delta a_\mu,\; \ldots \quad (\text{physics})
\end{align*}

All from $\varphi^2 = \varphi + 1$.

%==============================================================================
\section{Prerequisites}
%==============================================================================

\begin{enumerate}
\item Langlands--Tunnell (1981): icosahedral Artin $L$-functions are automorphic
\item Gelbart--Jacquet (1978): $\mathrm{Sym}^2$ lifting to $\mathrm{GL}(3)$
\item Jacquet--Shalika (1976): non-vanishing of $L(1, \mathrm{Sym}^2\rho)$
\item Weil positivity (unconditional): for admissible test functions, the zero sum is non-negative
\item Bochner's theorem: finite trigonometric polynomials have non-negative Fourier transform
\end{enumerate}

%==============================================================================
\section{The Shadow Theorem}
%==============================================================================

Irrational numbers are \textbf{shadows}---projections of rational 3D structure onto lower dimensions. The projection introduces irrationality. Squaring recovers the rational.

\begin{align*}
\sqrt{5}^2 &= 5 && \text{(1D shadow squared = integer)} \\
\pi^2 &\approx 2p = 10 && \text{(2D shadow squared = Schl\"afli base, 1.3\%)} \\
\pi^3 &\approx V + b_0 = 31 && \text{(2D shadow cubed = dodecahedral integer, 0.02\%)} \\
\pi^2 &= 2\!\left(p - \frac{\Delta}{\sqrt{5}}\right) && \text{(exact to 0.06\%)}
\end{align*}

The hierarchy:

\begin{itemize}[leftmargin=2em]
\item \textbf{1D shadows:} $\sqrt{5}$, $\varphi$ (irrational, algebraic)
\item \textbf{2D shadows:} $\pi$ (irrational, transcendental)
\item \textbf{3D structure:} $V=20$, $E=30$, $F=12$, $d=3$, $p=5$ (rational, integer, the source)
\end{itemize}

$\pi$ is not a constant. $\pi$ is what you get when you measure a 3D dodecahedron with a 2D ruler. The ruler cannot capture the full structure. The remainder = transcendence. The information lost in the projection. Square it: recover the rational structure.

Plato wrote in the \textit{Timaeus} (c.\ 360 BC): the cosmos is a dodecahedron. He wrote in the \textit{Republic}: reality is shadows on a wall. He had both pieces. He never connected them. The shadows on the wall ARE the irrational constants. The fire IS the axiom. The real objects ARE the dodecahedron. This connection: 2400 years later.

%==============================================================================
\section{Epilogue}
%==============================================================================

Pythagoras taught: all is number. His student Hippasus found $\sqrt{5}$ is irrational. They killed him for it.

The number they killed for generates the golden ratio, which generates the dodecahedron, which generates 137, which generates the critical line, which generates the Riemann Hypothesis, which generates the proton--electron mass ratio, which generates the cosmological constant.

The entire universe from a $1 \times 2$ rectangle and the theorem that bears his name.

All is number. He was right.

\begin{equation}
\boxed{\varphi^2 = \varphi + 1 \implies \text{RH}}
\end{equation}

\end{document}
